% ch29_monoidal_coherence.tex — Volume III Chapter 5

\chapter{Monoidal Categories: Coherence}

\begin{overviewbox}
Chapter~19 proved Mac Lane's coherence theorem for monoidal categories via
the normal-form strategy: every formal diagram of associators and unitors
in a free monoidal category collapses to a canonical bracketing, hence
commutes. Chapter~25 proved its bicategorical generalization, that every
bicategory is biequivalent to a strict 2-category.

This chapter is the synthesis. We re-derive monoidal coherence as a
\emph{special case} of bicategorical strictification — a monoidal category
is a one-object bicategory, so the strictification of Chapter~25 specializes
to the strictification of monoidal categories. We then push further into
material untouched by Chapter~19: \emph{coherence for monoidal functors}
(when does a monoidal functor's structure 2-cells satisfy ``all'' the
coherence diagrams automatically?), \emph{free monoidal categories} and
their universal property, and the symmetric/braided variants. The chapter
closes with a glimpse of operads — the algebraic objects that make
``monoidal coherence'' a workable computational tool.
\end{overviewbox}


% ═══════════════════════════════════════════════════════════════════════════
\section{Monoidal Strictification via Bicategorical Strictification}
\label{sec:mon-strict}
% ═══════════════════════════════════════════════════════════════════════════

\subsection{Formalizing}

Recall (Chapter~25, Proposition on monoidal $\Leftrightarrow$ one-object
bicategory): a monoidal category $(\mathcal{V}, \otimes, I, \alpha, \lambda,
\rho)$ is the same data as a one-object bicategory $B\mathcal{V}$.

\begin{theorem}[Strictification for monoidal categories]
\label{thm:mon-strict}
Every monoidal category $\mathcal{V}$ is monoidally equivalent to a
\emph{strict} monoidal category $\mathcal{V}^{\mathrm{st}}$.
\end{theorem}

\begin{prooflog}
\proofstep{Form the one-object bicategory $B\mathcal{V}$ (Chapter~25,
\S~\ref{sec:bicategories-formal}).}{Monoidal $\leftrightarrow$ one-object
bicategory.}
\proofstep{By Theorem~\ref{thm:strictification} (Chapter~25, bicategorical
strictification), $B\mathcal{V}$ is biequivalent to a strict 2-category
$(B\mathcal{V})^{\mathrm{st}}$.}{Apply the strictification construction.}
\proofstep{The string-replacement construction of Chapter~25 preserves the
one-object property: $(B\mathcal{V})^{\mathrm{st}}$ is again one-object.}{
Strings are constructed pointwise; the single 0-cell of $B\mathcal{V}$
remains the single 0-cell of $(B\mathcal{V})^{\mathrm{st}}$.}
\proofstep{Therefore $(B\mathcal{V})^{\mathrm{st}}$ corresponds to a strict
one-object 2-category, i.e., a strict monoidal category $\mathcal{V}^{\mathrm{st}}$.}{
Monoidal $\leftrightarrow$ one-object reverse direction.}
\proofstep{The biequivalence of bicategories $\widehat{(-)} :
(B\mathcal{V})^{\mathrm{st}} \to B\mathcal{V}$ corresponds to a monoidal
equivalence $\widehat{(-)} : \mathcal{V}^{\mathrm{st}} \to \mathcal{V}$.}{
A pseudofunctor between one-object bicategories is the same data as a
strong monoidal functor between the corresponding monoidal categories.}
\proofstep{Conclude: every monoidal category is monoidally equivalent to a
strict one.}{Strictification.}
\end{prooflog}

\begin{corollary}[Mac Lane's coherence, repackaged]
\label{cor:maclane-from-strict}
In any monoidal category, every diagram of associators and unitors commutes.
\end{corollary}

\begin{proof}
By Theorem~\ref{thm:mon-strict}, $\mathcal{V}$ is monoidally equivalent to
$\mathcal{V}^{\mathrm{st}}$, in which all $\alpha, \lambda, \rho$ are
identities and hence every formal diagram trivially commutes. The
equivalence reflects commutativity (it is fully faithful).
\end{proof}

\begin{remark}
Chapter~19's normal-form proof of coherence is constructive in $\mathcal{V}$
itself, while the strictification-based proof above factors through the
larger category $\mathcal{V}^{\mathrm{st}}$. Both yield the same theorem;
the strictification proof has the conceptual virtue of making coherence a
special case of the general bicategorical phenomenon.
\end{remark}

\subsection{Reorganizing: The Strictification Construction Explicitly}

The string-replacement construction of Chapter~25, specialized:

\begin{definition}[Explicit strictification $\mathcal{V}^{\mathrm{st}}$]
\label{def:explicit-strict}
For a monoidal category $\mathcal{V}$:
\begin{itemize}
  \item Objects of $\mathcal{V}^{\mathrm{st}}$: finite strings (lists) of
        objects of $\mathcal{V}$, including the empty string (which serves as
        the strict unit).
  \item Morphisms from $(A_1, \ldots, A_m)$ to $(B_1, \ldots, B_n)$: morphisms
        $A_1 \otimes \cdots \otimes A_m \to B_1 \otimes \cdots \otimes B_n$
        in $\mathcal{V}$ (left-associated), where the empty string composes
        to $I$.
  \item Tensor: concatenation of strings.
  \item Unit: empty string.
  \item Associator, unitors: identities.
\end{itemize}
\end{definition}

\begin{proposition}
$\mathcal{V}^{\mathrm{st}}$ is a strict monoidal category. The functor
$\widehat{(-)} : \mathcal{V}^{\mathrm{st}} \to \mathcal{V}$, $(A_1, \ldots,
A_n) \mapsto A_1 \otimes \cdots \otimes A_n$ (left-bracketed), is a monoidal
equivalence.
\end{proposition}

\begin{proof}
Concatenation is strictly associative; the empty string is a strict identity.
$\widehat{(-)}$ is essentially surjective (every $A \in \mathcal{V}$ is
$\widehat{(A)}$) and fully faithful (by definition of morphisms in
$\mathcal{V}^{\mathrm{st}}$).
\end{proof}


% ═══════════════════════════════════════════════════════════════════════════
\section{Coherence for Monoidal Functors}
\label{sec:monfun-coherence}
% ═══════════════════════════════════════════════════════════════════════════

A monoidal functor $F : \mathcal{V} \to \mathcal{W}$ between monoidal
categories carries comparison 2-cells $\phi_{A,B} : F A \otimes F B \to
F(A \otimes B)$ and $\phi_I : I_\mathcal{W} \to F I_\mathcal{V}$. The
coherence axioms in Chapter~19 ensure $\phi$ is compatible with the
associators and unitors of $\mathcal{V}, \mathcal{W}$.

\subsection{Formalizing}

\begin{theorem}[Coherence for monoidal functors]
\label{thm:monfun-coherence}
In a strong (resp.\ lax, oplax) monoidal functor $F : \mathcal{V} \to
\mathcal{W}$, every diagram built from $\phi_{A,B}$, $\phi_I$, $F$ applied to
the structure morphisms of $\mathcal{V}$, and the structure morphisms of
$\mathcal{W}$ — where the diagram has a single fixed source and target —
commutes.
\end{theorem}

\begin{remark}
The proof again reduces to strictification: replace $\mathcal{V}$ and
$\mathcal{W}$ by their strictifications, transport $F$ to a functor between
strict monoidal categories. In the strict case the coherence diagrams trivially
commute, and equivalence reflects.
\end{remark}

\subsection{Reorganizing: The Two Hexagons}

\begin{definition}[Hexagonal coherence axioms]
\label{def:monfun-hexagons}
A lax monoidal functor $F : (\mathcal{V}, \otimes, I) \to (\mathcal{W},
\otimes, I)$ satisfies the \term{associativity hexagon}: for $A, B, C \in
\mathcal{V}$,
\begin{equation*}
\begin{tikzcd}[column sep=tiny]
  (FA \otimes FB) \otimes FC \arrow[r, "{\alpha}"] \arrow[d, "{\phi_{A,B} \otimes 1}" '] &
  FA \otimes (FB \otimes FC) \arrow[d, "{1 \otimes \phi_{B,C}}"] \\
  F(A \otimes B) \otimes FC \arrow[d, "{\phi_{A \otimes B, C}}" '] &
  FA \otimes F(B \otimes C) \arrow[d, "{\phi_{A, B \otimes C}}"] \\
  F((A \otimes B) \otimes C) \arrow[r, "{F(\alpha)}"] & F(A \otimes (B \otimes C))
\end{tikzcd}
\end{equation*}
and the \term{unitality squares}:
\begin{equation*}
\begin{tikzcd}
  I_\mathcal{W} \otimes FA \arrow[r, "{\lambda_\mathcal{W}}"] \arrow[d, "{\phi_I \otimes 1}" '] &
  FA \\
  FI_\mathcal{V} \otimes FA \arrow[r, "{\phi_{I, A}}"] & F(I \otimes A) \arrow[u, "{F(\lambda)}" ']
\end{tikzcd}
\quad
\begin{tikzcd}
  FA \otimes I_\mathcal{W} \arrow[r, "{\rho_\mathcal{W}}"] \arrow[d, "{1 \otimes \phi_I}" '] &
  FA \\
  FA \otimes FI_\mathcal{V} \arrow[r, "{\phi_{A, I}}"] & F(A \otimes I) \arrow[u, "{F(\rho)}" ']
\end{tikzcd}
\end{equation*}
\end{definition}

\begin{proposition}[Coherence axioms imply all coherence diagrams]
\label{prop:hexagons-suffice}
The associativity hexagon and the unitality squares of
Definition~\ref{def:monfun-hexagons} imply Theorem~\ref{thm:monfun-coherence}:
all diagrams of formal coherence built from $\phi$, $F$, and the structure
2-cells of $\mathcal{V}, \mathcal{W}$ commute.
\end{proposition}

\begin{proof}[Sketch]
The strategy is the same as the monoidal-category case. Construct the
strictifications $\mathcal{V}^{\mathrm{st}}, \mathcal{W}^{\mathrm{st}}$;
transport $F$ to $F^{\mathrm{st}} : \mathcal{V}^{\mathrm{st}} \to
\mathcal{W}^{\mathrm{st}}$. In $\mathcal{W}^{\mathrm{st}}$, all coherence
diagrams have identity sides; the strictified $\phi^{\mathrm{st}}$ is
forced (by the two hexagons and squares) to be an identity. Pulling back,
all original diagrams commute.
\end{proof}


% ═══════════════════════════════════════════════════════════════════════════
\section{Free Monoidal Categories}
\label{sec:free-monoidal}
% ═══════════════════════════════════════════════════════════════════════════

\subsection{Formalizing}

The universal monoidal category on a given category $\mathcal{C}$ is the
``one with no relations beyond what monoidal structure requires.'' Formally:

\begin{definition}[Free monoidal category]
\label{def:free-monoidal}
For an ordinary category $\mathcal{C}$, the \term{free monoidal category} on
$\mathcal{C}$, written $F^\otimes(\mathcal{C})$, is the monoidal category
characterized by the universal property: for every monoidal category
$\mathcal{V}$ and every ordinary functor $F : \mathcal{C} \to \mathcal{V}$,
there is an essentially unique strong monoidal functor $\widetilde F :
F^\otimes(\mathcal{C}) \to \mathcal{V}$ extending $F$.
\end{definition}

\begin{theorem}[Construction]
\label{thm:free-monoidal-construction}
$F^\otimes(\mathcal{C})$ has:
\begin{itemize}
  \item Objects: finite lists (with bracketings, possibly nested) of objects of $\mathcal{C}$.
  \item Morphisms: generated by morphisms of $\mathcal{C}$ on singletons,
        identity for tensor combinations, with associators and unitors
        formally adjoined as isomorphisms.
\end{itemize}
The universal property follows from a normal-form argument akin to
strictification.
\end{theorem}

\begin{proof}[Sketch]
Existence: the explicit construction. Universality: given $F : \mathcal{C}
\to \mathcal{V}$, define $\widetilde F$ on singletons by $F$; extend to
tensor combinations by $F(A_1) \otimes \cdots \otimes F(A_n)$ (with
appropriate bracketing) and on morphisms by transport along the
$\mathcal{V}$-associators and unitors. Coherence (Theorem~\ref{thm:monfun-coherence})
guarantees the extension is well-defined and a strong monoidal functor.
Uniqueness up to monoidal natural isomorphism is standard.
\end{proof}

\begin{remark}[Strict variant]
The \emph{strict} free monoidal category $F^{\otimes,\mathrm{st}}(\mathcal{C})$
has objects flat lists (no nested brackets) and morphisms generated by $\mathcal{C}$
without formal associators. It is the strict 2-category one-object specialization
of the strictification construction.
\end{remark}

\subsection{Reorganizing: The Yoneda Description}

\begin{proposition}[Free monoidal as Day convolution]
\label{prop:free-mon-day}
For a category $\mathcal{C}$ (with no monoidal structure), $F^\otimes(\mathcal{C})$
is the presheaf category $[\mathcal{C}^{\mathrm{op}}, \mathbf{Set}]$ equipped
with \emph{Day convolution} — see Chapter~30. The Yoneda embedding $Y :
\mathcal{C} \to [\mathcal{C}^{\mathrm{op}}, \mathbf{Set}]$ becomes the
unit of $F^\otimes$.
\end{proposition}

This identification — free cocompletion = free monoidal category, in the
appropriate sense — is the link between Yoneda theory (Chapter~28) and
monoidal theory developed in Chapter~30.


% ═══════════════════════════════════════════════════════════════════════════
\section{Symmetric and Braided Coherence}
\label{sec:sym-coherence}
% ═══════════════════════════════════════════════════════════════════════════

Symmetric and braided monoidal categories add an extra structure 2-cell
$\sigma_{A,B} : A \otimes B \cong B \otimes A$ (the \term{braiding} or
\term{symmetry}), satisfying additional hexagons.

\subsection{Formalizing}

\begin{definition}[Braided monoidal category]
\label{def:braided}
A \term{braided monoidal category} adds to a monoidal category a natural
isomorphism $\sigma_{A, B} : A \otimes B \to B \otimes A$ — the \term{braiding}
— satisfying the two \term{braiding hexagons} (compatibility with the
associator).
\end{definition}

\begin{definition}[Symmetric monoidal category]
A \term{symmetric monoidal category} is a braided monoidal category whose
braiding is \term{self-inverse}: $\sigma_{B, A} \circ \sigma_{A, B} =
\mathrm{id}_{A \otimes B}$.
\end{definition}

\subsection{Formally}

\begin{theorem}[Coherence for symmetric monoidal categories]
\label{thm:sym-coherence}
In a symmetric monoidal category, every diagram built from $\alpha, \lambda,
\rho, \sigma$ between two parenthesizations of the same word — \emph{where
both words use the same permutation of the input letters} — commutes.
\end{theorem}

\begin{remark}[Note the constraint]
Theorem~\ref{thm:sym-coherence} is \emph{not} ``every formal diagram
commutes.'' Two different ways of permuting $A \otimes B \otimes C$ via
braidings can give two different morphisms $A \otimes B \otimes C \to
B \otimes A \otimes C$. The theorem says: if both ends of the diagram
correspond to the \emph{same} permutation, then any path connecting them
commutes. Braids are real morphisms, not coherence noise.
\end{remark}

\begin{theorem}[Coherence for braided monoidal categories]
\label{thm:braid-coherence}
In a braided monoidal category, the coherence ``free model'' is the braid
category $\mathbf{Br}$: objects are natural numbers, morphisms from $m$ to
$m$ are braids on $m$ strands. Two formal diagrams commute iff the
corresponding braids are equal as braids.
\end{theorem}

\begin{remark}
The contrast: symmetric coherence quotients out the over-and-under information
in braids, treating each as a permutation; braided coherence keeps the full
braid. This is why braided monoidal categories model topological information
(knots, links, ribbon categories) that symmetric ones miss.
\end{remark}


% ═══════════════════════════════════════════════════════════════════════════
\section{Operads: A Glimpse}
\label{sec:operads}
% ═══════════════════════════════════════════════════════════════════════════

The combinatorial structure underlying monoidal coherence is captured by
\term{operads} (and their relatives — props, properads, multicategories).

\subsection{Formalizing}

\begin{definition}[Operad]
\label{def:operad}
A \term{(symmetric, $\mathcal{V}$-enriched) operad} $\mathcal{O}$ consists of:
\begin{itemize}
  \item For each $n \geq 0$, an object $\mathcal{O}(n) \in \mathcal{V}$
        equipped with a right action of the symmetric group $\mathfrak{S}_n$.
  \item Composition maps $\gamma : \mathcal{O}(n) \otimes \mathcal{O}(k_1)
        \otimes \cdots \otimes \mathcal{O}(k_n) \to \mathcal{O}(k_1 + \cdots
        + k_n)$ for all $n, k_1, \ldots, k_n$.
  \item Unit $\eta : I \to \mathcal{O}(1)$.
\end{itemize}
Satisfying associativity, unit, and equivariance axioms.
\end{definition}

\begin{remark}[Algebras over an operad]
For an operad $\mathcal{O}$ in $\mathcal{V}$, an \term{$\mathcal{O}$-algebra}
is an object $A \in \mathcal{V}$ equipped with structure morphisms
$\mathcal{O}(n) \otimes A^{\otimes n} \to A$ satisfying compatibility with
$\gamma$ and $\eta$.

\emph{Examples:}
\begin{itemize}
  \item The \emph{associative operad} $\mathrm{Assoc}$ in $\mathbf{Set}$:
        $\mathrm{Assoc}(n) = \mathfrak{S}_n$. Algebras are non-commutative monoids.
  \item The \emph{commutative operad} $\mathrm{Comm}$: $\mathrm{Comm}(n) =
        \{*\}$. Algebras are commutative monoids.
  \item The \emph{little disks operad} (in $\mathbf{Top}$): $E_2$-algebras
        are braided monoidal structures up to homotopy.
\end{itemize}
\end{remark}

\subsection{Reorganizing: Operads as Monads}

\begin{proposition}[Operad-monad correspondence]
\label{prop:operad-monad}
A symmetric $\mathcal{V}$-operad $\mathcal{O}$ in a cocomplete monoidal
$\mathcal{V}$ gives a monad $T_\mathcal{O}$ on $\mathcal{V}$:
\begin{equation*}
  T_\mathcal{O}(X) \;=\; \coprod_{n \geq 0} \mathcal{O}(n) \otimes_{\mathfrak{S}_n}
  X^{\otimes n}.
\end{equation*}
$\mathcal{O}$-algebras and $T_\mathcal{O}$-algebras coincide.
\end{proposition}

\begin{remark}
This is the link between operad theory and the formal theory of monads
(Chapter~32). An operad presents an algebraic theory with $n$-ary operations
for every $n$; the monad presents the same theory as a coalgebraic gadget.
\end{remark}


% ═══════════════════════════════════════════════════════════════════════════
\section{Exercises}
\label{sec:mon-coherence-exercises}
% ═══════════════════════════════════════════════════════════════════════════

\begin{exercisesbox}
\begin{qa}
\qaitem{State the monoidal strictification theorem.}{Every monoidal category is monoidally equivalent to a strict monoidal category.}
\qaitem{How does monoidal strictification follow from bicategorical strictification?}{A monoidal category is a one-object bicategory. The bicategorical strictification of Chapter~25 preserves the one-object property; the resulting strict 2-category corresponds to a strict monoidal category.}
\qaitem{Describe $\mathcal{V}^{\mathrm{st}}$ explicitly.}{Objects: finite lists of objects of $\mathcal{V}$. Morphisms: morphisms of $\mathcal{V}$ between the left-associated tensor of the lists. Tensor: concatenation. Unit: empty list. All structure 2-cells are identities.}
\qaitem{What is the equivalence $\widehat{(-)} : \mathcal{V}^{\mathrm{st}} \to \mathcal{V}$?}{$(A_1, \ldots, A_n) \mapsto A_1 \otimes \cdots \otimes A_n$ (left-bracketed). It is essentially surjective and fully faithful.}
\qaitem{State Mac Lane's coherence theorem (in this language).}{In any monoidal category, every diagram of associators and unitors with the same source and target commutes.}
\qaitem{Why does Mac Lane's coherence follow from strictification?}{In the strict category, all diagrams of identity 2-cells trivially commute. The equivalence with $\mathcal{V}$ reflects this.}
\qaitem{State coherence for monoidal functors.}{In a (lax, oplax, or strong) monoidal functor $F : \mathcal{V} \to \mathcal{W}$, every diagram of formal $F$-coherence morphisms (involving $\phi_{A,B}$, $\phi_I$, $F$ applied to $\mathcal{V}$'s structure morphisms, and $\mathcal{W}$'s structure morphisms) commutes.}
\qaitem{What axioms guarantee monoidal-functor coherence?}{The associativity hexagon and the unitality squares of Definition~\ref{def:monfun-hexagons}.}
\qaitem{Define the free monoidal category $F^\otimes(\mathcal{C})$.}{The monoidal category with the universal property: every functor $\mathcal{C} \to \mathcal{V}$ extends uniquely (up to monoidal natural iso) to a strong monoidal functor $F^\otimes(\mathcal{C}) \to \mathcal{V}$.}
\qaitem{Describe $F^\otimes(\mathcal{C})$ explicitly.}{Objects: finite bracketed lists of objects of $\mathcal{C}$. Morphisms: generated by morphisms of $\mathcal{C}$ on singletons and the associators/unitors as formal isomorphisms.}
\qaitem{What is the relationship between $F^\otimes(\mathcal{C})$ and the presheaf category?}{$F^\otimes(\mathcal{C})$ is the presheaf category $[\mathcal{C}^{\mathrm{op}}, \mathbf{Set}]$ with Day convolution as the tensor (Chapter~30).}
\qaitem{Define a braided monoidal category.}{A monoidal category with a natural isomorphism $\sigma_{A,B} : A \otimes B \to B \otimes A$ — the braiding — satisfying two compatibility hexagons with the associator.}
\qaitem{Define a symmetric monoidal category.}{A braided monoidal category whose braiding is self-inverse: $\sigma_{B,A} \circ \sigma_{A,B} = \mathrm{id}$.}
\qaitem{State coherence for symmetric monoidal categories.}{Every diagram of $\alpha, \lambda, \rho, \sigma$ between two parenthesizations of the same word — using the same permutation — commutes.}
\qaitem{Why is symmetric coherence not ``every diagram commutes''?}{Because different choices of braiding can give genuinely different permutations; coherence quotients by permutation equivalence, not all rearrangements.}
\qaitem{State coherence for braided monoidal categories.}{The coherence-free model is the braid category $\mathbf{Br}$. Two formal diagrams commute iff their corresponding braids are equal as braids.}
\qaitem{What is the difference between symmetric and braided coherence?}{Symmetric coherence forgets the ``over-vs.\ under'' information in a braid (keeps only the permutation). Braided coherence retains the full braid. Hence braided monoidal categories model knots/links/ribbon structures while symmetric ones don't.}
\qaitem{Define an operad.}{A family $\{\mathcal{O}(n)\}_{n \geq 0}$ of $\mathfrak{S}_n$-equipped objects with composition maps $\gamma : \mathcal{O}(n) \otimes \prod_i \mathcal{O}(k_i) \to \mathcal{O}(\sum k_i)$ and unit, satisfying associativity, unit, and equivariance.}
\qaitem{What is an algebra over an operad $\mathcal{O}$?}{An object $A$ with structure maps $\mathcal{O}(n) \otimes A^{\otimes n} \to A$ compatible with $\gamma$ and $\eta$.}
\qaitem{Give three example operads.}{$\mathrm{Assoc}$ (associative operad, algebras = monoids); $\mathrm{Comm}$ (commutative operad, algebras = commutative monoids); $E_2$ (little disks operad in $\mathbf{Top}$, algebras = braided monoidal structures up to homotopy).}
\qaitem{State the operad-monad correspondence.}{Each operad $\mathcal{O}$ gives a monad $T_\mathcal{O}(X) = \coprod_n \mathcal{O}(n) \otimes_{\mathfrak{S}_n} X^{\otimes n}$; algebras coincide.}
\qaitem{Why do operads matter for coherence?}{They make the combinatorial structure of monoidal coherence (and braided/symmetric variants) into a workable algebraic gadget. Operad theory governs which coherence diagrams must commute.}
\qaitem{What is the strict free monoidal category $F^{\otimes,\mathrm{st}}(\mathcal{C})$?}{Objects: flat lists. Morphisms: generated by $\mathcal{C}$ alone, no formal associators. It is the strict 2-category specialization of strictification.}
\qaitem{How does the strictification of monoidal functors work?}{Given $F : \mathcal{V} \to \mathcal{W}$, transport to $F^{\mathrm{st}} : \mathcal{V}^{\mathrm{st}} \to \mathcal{W}^{\mathrm{st}}$. Since $\mathcal{W}^{\mathrm{st}}$ is strict, $F^{\mathrm{st}}$'s coherence comparators must reduce; the axioms force them to be identities.}
\qaitem{In what sense is coherence ``conceptual'' rather than ``combinatorial''?}{Once strictification is in hand, coherence reduces to the trivial fact that diagrams in strict categories have identity sides. The strictification theorem absorbs all combinatorial intricacy.}
\qaitem{Why is the bicategorical strictification proof of Chapter~25 ``deeper'' than the monoidal case of Chapter~19?}{Because Chapter~19 used a normal-form combinatorial argument specialized to monoidal coherence. Chapter~25's strictification works for arbitrary bicategories — and monoidal is just one example. The bicategorical proof unifies many coherence theorems.}
\qaitem{What is the ``free strict monoidal category on a strictly monoidal generating data''?}{The free strict monoidal category $F^{\otimes,\mathrm{st}}(\mathcal{C})$ is its own strictification.}
\qaitem{What is the monoidal natural transformation between two strong monoidal functors?}{A natural transformation $\eta : F \Rightarrow G$ such that the diagrams $\eta_{A \otimes B} \circ \phi^F_{A,B} = \phi^G_{A,B} \circ (\eta_A \otimes \eta_B)$ and $\eta_I \circ \phi^F_I = \phi^G_I$ commute.}
\qaitem{What does monoidal equivalence buy beyond monoidal isomorphism?}{Monoidal equivalence is weaker: it requires only essentially surjective fully faithful strong monoidal functors with monoidal natural isomorphisms as units/counits. Strictification gives \emph{equivalence}, not isomorphism, between $\mathcal{V}$ and $\mathcal{V}^{\mathrm{st}}$.}
\qaitem{Summarize the chapter's contribution beyond Chapter~19.}{Coherence is now a bicategorical phenomenon (one-object special case of strictification); monoidal-functor coherence is established formally; free monoidal categories and operads give the algebraic backbone.}
\end{qa}
\end{exercisesbox}


% ═══════════════════════════════════════════════════════════════════════════
\section*{Chapter Summary}
\addcontentsline{toc}{section}{Chapter Summary}
% ═══════════════════════════════════════════════════════════════════════════

\begin{summarybox}
\textbf{Strictification, derived twice.}\quad Every monoidal category is
monoidally equivalent to a strict one — once proved combinatorially in
Chapter~19, here re-proved as the one-object special case of bicategorical
strictification (Chapter~25). Mac Lane's coherence becomes the trivial
observation that in the strict category, all diagrams have identity sides.

\textbf{Coherence for monoidal functors.}\quad The associativity hexagon and
unitality squares of $\phi_{A,B}, \phi_I$ imply that every formal coherence
diagram involving $F$, its comparators, and the structure of source/target
commutes. Same strictification proof.

\textbf{Free monoidal categories.}\quad $F^\otimes(\mathcal{C})$ — objects =
bracketed lists, morphisms generated by $\mathcal{C}$ and formal coherence
isos — is universal among monoidal categories receiving a functor from
$\mathcal{C}$. It equals the presheaf category with Day convolution
(Chapter~30).

\textbf{Symmetric/braided coherence.}\quad Symmetric coherence: every diagram
between two parenthesizations of the same word using the same permutation
commutes (different permutations are genuinely different morphisms). Braided
coherence: governed by the braid category $\mathbf{Br}$. The distinction
explains why braided monoidal categories model knots and ribbons.

\textbf{Operads.}\quad The algebraic gadget capturing $n$-ary structure
under monoidal coherence. Every operad yields a monad; algebras coincide.
Examples: $\mathrm{Assoc}$, $\mathrm{Comm}$, $E_2$. Operads are the
combinatorial spine of higher coherence theory.
\end{summarybox}


% ═══════════════════════════════════════════════════════════════════════════
\section*{Formal Restatement}
\addcontentsline{toc}{section}{Formal Restatement}
% ═══════════════════════════════════════════════════════════════════════════

\begin{formalrestatement}
\textbf{Strictification.}\quad For every monoidal category $\mathcal{V}$,
there exist a strict monoidal category $\mathcal{V}^{\mathrm{st}}$ and a
strong monoidal equivalence $\widehat{(-)} : \mathcal{V}^{\mathrm{st}} \to
\mathcal{V}$. Construction: objects are finite lists of objects of
$\mathcal{V}$; morphisms are $\mathcal{V}$-morphisms between left-bracketed
tensors; tensor is concatenation.

\medskip
\textbf{Mac Lane coherence.}\quad Every diagram of associators and unitors
in a monoidal category commutes (special case of strictification).

\medskip
\textbf{Monoidal-functor coherence.}\quad If $F : \mathcal{V} \to \mathcal{W}$
is monoidal (lax/oplax/strong) satisfying the associativity hexagon and
unitality squares, every diagram of formal coherence commutes.

\medskip
\textbf{Free monoidal category.}\quad For any category $\mathcal{C}$,
$F^\otimes(\mathcal{C})$ exists and satisfies the universal property: every
functor $\mathcal{C} \to \mathcal{V}$ extends essentially uniquely to a
strong monoidal functor $F^\otimes(\mathcal{C}) \to \mathcal{V}$.

\medskip
\textbf{Symmetric coherence.}\quad In a symmetric monoidal category, every
diagram of $\alpha, \lambda, \rho, \sigma$ between parenthesizations
representing the \emph{same permutation} of the input letters commutes.

\medskip
\textbf{Braided coherence.}\quad The coherence-free model is the braid
category $\mathbf{Br}$; formal diagrams commute iff their braids are equal.

\medskip
\textbf{Operad to monad.}\quad For a symmetric $\mathcal{V}$-operad
$\mathcal{O}$ in cocomplete $\mathcal{V}$:
\begin{equation*}
  T_\mathcal{O}(X) \;=\; \coprod_{n \geq 0} \mathcal{O}(n) \otimes_{\mathfrak{S}_n}
  X^{\otimes n},
\end{equation*}
with $\mathcal{O}$-algebras = $T_\mathcal{O}$-algebras.

\medskip
\noindent\textit{Chapter~30 introduces \emph{Day convolution} — the
monoidal structure on enriched functor categories that completes the
identification of $F^\otimes(\mathcal{C})$ with $[\mathcal{C}^{\mathrm{op}},
\mathbf{Set}]$.}
\end{formalrestatement}
